\scanpage{580}

% source-position: upper portion, presentations (30)--(31).
In terms of $\rho,\sigma,\tau'$ we find likewise
\[
(30)\qquad
\left\{
\begin{aligned}
&\rho^6=\sigma^4=\tau'^2=1,\quad \rho^3=\sigma^2\\
&\tau'(\rho)=\rho^{-1},\quad\tau'(\sigma)=\sigma^{-1}\\
&\sigma(\rho)=\rho^{-1}
\end{aligned}
\right.
\]
The subgroup $\mathcal G^+/\Pi_0$ has the presentation in $\rho,\sigma$
\[
(31)\qquad
\left\{
\begin{aligned}
&\rho^6=\sigma^4=1,\quad \rho^3=\sigma^2\\
&\sigma(\rho)=\rho^{-1}
\end{aligned}
\right\}
\]

% source-position: upper-middle portion, action question and extension diagram (32).
What is the operation of $\mathcal M_{0,3}^{\sim}\simeq\mathcal M/\pm1$
on $\mathcal G^+/\Pi_0$ ?
\deleted{\unsure{it is an} operation}
\inserted{We have a structure} of extension
\[
(32)\qquad
\diagram{
\begin{tikzcd}
1 \arrow[r] & \mu \arrow[r] \arrow[d,equal]
& \mathcal G^+/\Pi_0 \arrow[r]
& \mathcal G^+/\Pi \arrow[r] \arrow[d,equal]
& 1\\
& \pm1 && \mathcal G_3 &
\end{tikzcd}
}
\]
% source-position: middle portion, prose following (32).
and the \unsure{automorphisms} in question preserve this extension structure.
The \unsure{elements|elements} of
$\mathcal M_{0,3}^{!\sim}\simeq\mathcal M^!/\{\pm1\}$
operate by inducing the identity in the quotient $\mathcal G_3$,
and of \unsure{course} also in $\mu$ --- now such an automorphism of
the extension corresponds to a homomorphism
% source-position: lower-middle portion, character maps (33).
\[
(33)\qquad
\diagram{
\begin{tikzcd}
\mathcal G_3 \arrow[r] & \mu \arrow[d,equal]\\
& \pm1
\end{tikzcd}
}
\quad\text{i.e.}\quad
\diagram{
\begin{tikzcd}
\mathcal G_{3\,\mathrm{ab}} \arrow[r] & \mu
\end{tikzcd}
}
\]
% source-position: lower portion, sign character and conclusion (34).
now
\[
\diagram{
\begin{tikzcd}
\mathcal G_{3\,\mathrm{ab}}\arrow[r,"\uline{\mathrm{sg}}"]
&\{\pm1\}=\mu
\end{tikzcd}
}
\]
thus the group of automorphisms of $\mathcal G^+/\Pi_0$ inducing
the identity in $\mathcal G_3$ is $\simeq\mathbb Z/2\mathbb Z$,
a generator \inserted{$\Theta$} is given by
\[
(34)\qquad
\Theta(x)=\operatorname{sg}(x)\mathbin{.}x
\quad\text{thus}\quad
\left\{
\begin{aligned}
\Theta(\rho)&=\rho\\
\Theta(\sigma)&=-\sigma
\end{aligned}
\right.
\]

\mspage{579}

% Printed archive footer outside the manuscript frame (not manuscript):
% 579
% https://grothendieck.umontpellier.fr
% Université de Montpellier, UMR CNRS 5149, Case courrier 051, Place Eugène Bataillon
% 34095 Montpellier cedex 5 - France

\scanpage{581}
\mspage{594}

% source-position: upper portion, displays (35) and (35 bis).
Comparing with the formulas $(20)$, we find
\[
(35)\qquad
\deleted{\unsure{u}}
\left.u\right|_{\mathcal G^{+\wedge}/\widehat{\Pi}_0}
=\Theta^{\varepsilon_2(u)}:
\quad
\diagram{
\begin{tikzcd}
\dot\rho\arrow[r,mapsto]&\dot\rho\\
\dot\sigma\arrow[r,mapsto]&\varepsilon_2(u)\dot\sigma
\end{tikzcd}
}
\]
i.e.
\[
(35\text{ bis})\qquad
\begin{gathered}
[\dot u,\dot\sigma]\\[-0.5ex]
\unsure{\Vert}\\[-0.5ex]
\inserted{\dot u(\dot\sigma)\dot\sigma^{-1}}
\end{gathered}
=\varepsilon_2(u)^{\cdot}
\quad\text{in }\mathcal G/\Pi_0.
\]
% source-position: upper left slanted margin, commutator formula and warning.
\margin{This can also be written
$[\dot u,\dot x]=\operatorname{sg}(x)^{\unsure{\nu_2(u)}}$
in $\mathcal G^{+\wedge}/\Pi_0$
$(u\in\mathcal M^!,x\in\mathcal G^{+\wedge})$
(careful, this \unsure{goes wrong|does} \unsure{if we assume there}
\unsure{for} $u\in\mathcal M^!$ only $u\in\mathcal M\ldots$)}

% source-position: upper-middle portion, revised heading and definition (36).
\deleted{I would like}
Let us take up again the definition
\[
(36)\qquad
Z=Z_{\rho,\sigma}\stackrel{\mathrm{dfn}}=
\left\{
u\in\operatorname{Aut}(\mathcal G^{+\wedge})\ \left|\
\begin{aligned}
&u(\unsure{L_\rho})\
\inserted{\mathcal G^{+\wedge}\text{-}}\text{conjugate to }\
\deleted{\unsure{\ell_\rho}}L_\rho\\
&u(\unsure{L_\sigma})\
\inserted{\mathcal G^{+\wedge}}\text{-conjugate : }L_\sigma\\
&u(\unsure{L_0})=L_0
\end{aligned}
\right.
\right\}
\]
so that we shall have, for $u\in Z$
% source-position: middle portion, display (37) and parameter conditions.
\[
(37)\qquad
\left\{
\begin{aligned}
&u(\rho)\quad\mathcal G\text{-conjugate to }\rho^{\varepsilon'}\\
&u(\sigma)\quad\mathcal G\text{-conjugate to }\sigma^{\unsure{\varepsilon}}
\quad(=\varepsilon\sigma)\\
&u(\varepsilon_0)=\varepsilon_0^\mu
\end{aligned}
\right.
\]
with $\mu\in\widehat{\mathbb Z}^{*}$ and
$\varepsilon,\varepsilon'\in\{\pm1\}$ well determined,
(\unsure{for}
\diagram{
\begin{tikzcd}
\widehat{\mathbb Z}\arrow[r]&\mathcal G\\
n\arrow[r,mapsto]&\varepsilon_0^n
\end{tikzcd}
}
is injective, $\rho$ and $\rho^{-1}$ are not
$\mathcal G^{+\wedge}$-conjugate, nor $\sigma$ and $\sigma^{-1}$).
% source-position: lower-middle portion, proof of stabilization and l_0, l_1 calculations.
I say that such a $u$ comes from $\mathcal M_{0,3}^{\sim}$ ---
it remains to check, by virtue of the proposition..., that $u$ stabilizes
$\widehat{\Pi}$, which is generated by $\ell_0,\ell_1,\ell_\infty$.
Now we shall have
\[
u(\ell_0)=\ell_0^\mu\in\widehat{\Pi}
\]
and
\[
u(\ell_1)=
\underbrace{u\rho}_{\inserted{
\begin{gathered}
\deleted{\unsure{\operatorname{int}(f)\rho u}}\\
\deleted{\unsure{\operatorname{int}(f)\rho^{\varepsilon'}u}}\\
\operatorname{int}(f)(\rho^{\varepsilon'})u\\
f\in\mathcal G^{+\wedge}
\end{gathered}
}}(\ell_0)
=\operatorname{int}(f)(\rho^{\varepsilon'})(u(\ell_0))
=\operatorname{int}(f)(\rho)^{\varepsilon'}(\ell_0^\mu)
\]
now $\mathcal G^\wedge$ and a fortiori
$\operatorname{int}(f)(\rho)^{\varepsilon'}$ normalizes
$\deleted{\unsure{\mathcal G}}\widehat{\Pi}$,
thus \deleted{and} $u(\ell_1)\in\widehat{\Pi}$
% source-position: lower portion, l_infinity calculation and final cancelled words.
and since $u(-1)=-1$,
\inserted{whence:}
\[
u(\ell_\infty)=-u(\ell_1\ell_0)^{-1}
=-u(\ell_0^{-1})u(\ell_1)^{-1}
\]
we also have $u(\ell_\infty)\in\Pi$, OK.

If we identify $u$ with an element of $\mathcal M_{0,3}^{\sim}$,
we see that
\deleted{\unsure{it is}} in fact
\deleted{an element of $\unsure{Z|\mathcal M}$.}

\mspage{580}

% Printed archive footer outside the manuscript frame (not manuscript):
% 580
% https://grothendieck.umontpellier.fr
% Université de Montpellier, UMR CNRS 5149, Case courrier 051, Place Eugène Bataillon
% 34095 Montpellier cedex 5 - France

\scanpage{582}

% source-position: upper portion, continuation from preceding page.
\[
u\in\mathcal M_{0,3}^{\sim}[0]\simeq\mathcal M^![0]\subset\mathcal M
\]
% source-position: large slanted upper-left margin, correction and later-page reference.
\margin{Careful
we had not previously introduced
$\unsure{\mathcal M_{0,3}^{\sim}|\mathcal M_{0,3}}[0]$, but only
$\mathcal M_{0,3}^{!\sim}[0]\subset\mathcal M_{0,3}^{!\sim}$.
Now it \unsure{is not}
true that we shall have here $u\in\mathcal M_{0,3}^{!\sim}$,
for example \unsure{we may have} $u=\varepsilon_0$.
The image of $u$ in $\mathcal G_3$ is in
\deleted{\unsure{\operatorname{sg}}}
$\{1,\dot\sigma_0\}$, thus by correcting
\inserted{by $\unsure{\varepsilon_0}$},
we shall have an element of $\mathcal M_{0,3}^{!\sim}[0]$.
\deleted{We must}
Cf. more detailed correction on page $609$ ff.
$Z_{\rho,\sigma}^!=\mathcal M^![0]=Z_{\rho,\sigma}\cap\mathcal M^!$.}

% source-position: upper portion, display (38) and its cancelled side conditions.
\deleted{\unsure{we have}} thus $(20)$, more precisely
\[
(38)\qquad
\left\{
\begin{aligned}
&u(\rho)=\operatorname{int}(\alpha)\rho\\
&u(\sigma)=\varepsilon_2\operatorname{int}(\beta)\sigma\\
&u(\varepsilon_0)=\varepsilon_0^\mu
\end{aligned}
\right.
\qquad
\begin{aligned}
&\unsure{\text{with}}\ \deleted{\unsure{\alpha}}
\quad\alpha\in\widehat{\Pi}_{0\tau},\ \beta\in\widehat{\Pi}_0\\
&\deleted{\unsure{\varepsilon_2}}\quad
\left\{\begin{aligned}
&\mu\in\widehat{\mathbb Z}^{*}\\
&\deleted{\varepsilon_2=\varepsilon_2(\mu),\
\varepsilon_\tau(\alpha)=\varepsilon_3(\mu)}
\end{aligned}\right.
\end{aligned}
\]
% source-position: upper-middle portion, uniqueness paragraph after (38).
\unsure{in} their turn \unsure{imply|implying} $(37)$. In fact
the formulas \inserted{$(38)$} for $u(\rho),u(\sigma),u(\varepsilon_0)$
uniquely determine
$\mu\in\widehat{\mathbb Z}^{*}$, $\varepsilon_2\in\{\pm1\}$,
$\alpha\in\widehat{\Pi}_{0,\tau}$, $\beta\in\widehat{\Pi}_0$,
and we shall then necessarily have
% source-position: middle portion, display (39), exact sequence (40).
\[
(39)\qquad
\varepsilon_2=\varepsilon_2(\mu),\qquad
\varepsilon_\tau(\alpha)=\varepsilon_3(\mu)
\]
where $\varepsilon_\tau$ is the homomorphism that occurs in the
\inserted{exact} sequence
\[
(40)\qquad
\deleted{\unsure{\varepsilon_0}}
\diagram{
\begin{tikzcd}
1\arrow[r]&\widehat{\Pi}_0\arrow[r]&
\widehat{\Pi}_{0,\tau}\arrow[r,"\varepsilon_\tau"]&
\{\pm1\}\arrow[r]&1
\end{tikzcd}
}.
\]

% source-position: lower-middle portion, proposition and component identifications.
Thus we have proved

\uline{Proposition}\quad
In the universal case $(\mathcal G=\mathcal G^{+\wedge})$
we have
\[
\diagram{
\begin{tikzcd}
Z_{\rho,\sigma}^!&
\mathcal M^![0]\stackrel{\mathrm{dfn}}=
\mathcal M_{0,3}^{!\sim}[0]\arrow[l,"\sim"']
\end{tikzcd}
},
\]
the components $u,\mu,\varepsilon',\varepsilon$ of an element of
$Z_{\rho,\sigma}^!$ correspond resp.
(in terms of $u_0\in\mathcal M[0]$):
$u_0$ identified with an element of
$\operatorname{Aut}(\mathcal G^{+\wedge})$,
$\mu=\chi(u_0)$ (the multiplier of $u_0$),
$\varepsilon=\varepsilon_3(\mu)$,
$\varepsilon'=\varepsilon_2(\mu)$.

% source-position: lower portion, corollary and chain of quotients.
\uline{Corollary}\quad We have
\[
G_{\rho,\sigma}\simeq\mathcal M=\mathcal N_{1,1}^{\sim},
\]
\[
\gamma_{\rho,\sigma}\simeq
\mathcal M/\mathcal G^{+\wedge}
\simeq\mathcal N_{1,1}^{\sim}/\mathcal G^{+\wedge}
\simeq\mathcal M_{0,3}^{\sim}/\mathcal G_{0,3}^{+\wedge}
\simeq\Gamma_{\mathbb Q}^{\sim}.
\]

\mspage{581}

% Printed archive footer outside the manuscript frame (not manuscript):
% 581
% https://grothendieck.umontpellier.fr
% Université de Montpellier, UMR CNRS 5149, Case courrier 051, Place Eugène Bataillon
% 34095 Montpellier cedex 5 - France